% SPDX-License-Identifier: MIT
% Copyright (c) 2026 Evgeny Baulin
% Optimization Methods, Seminar 06: one-page cheat sheet (Russian).
% Generated by checks/06. Second-Order Methods Part Two.ipynb from web/06/js/content-cheatsheet.js and content/students.json. Do not edit by hand.
% Compile with: latexmk -xelatex -interaction=nonstopmode -halt-on-error Cheatsheet_ru.tex
\RequirePackage{fix-cm}
\documentclass[10pt,a4paper]{article}

\usepackage{amsmath,amssymb}
\usepackage{fontspec}
\setmainfont{cmunrm.otf}[BoldFont=cmunbx.otf, ItalicFont=cmunti.otf, BoldItalicFont=cmunbi.otf]
\setsansfont{cmunss.otf}[BoldFont=cmunsx.otf, ItalicFont=cmunsi.otf, BoldItalicFont=cmunso.otf]
\setmonofont{cmuntt.otf}[BoldFont=cmuntb.otf, ItalicFont=cmunit.otf, BoldItalicFont=cmuntx.otf]
\usepackage{polyglossia}
\setdefaultlanguage{russian}
\setotherlanguage{english}
\newfontfamily\cyrillicfont{cmunrm.otf}[Ligatures=TeX, BoldFont=cmunbx.otf, ItalicFont=cmunti.otf, BoldItalicFont=cmunbi.otf]
\newfontfamily\cyrillicfontsf{cmunss.otf}[Ligatures=TeX, BoldFont=cmunsx.otf, ItalicFont=cmunsi.otf, BoldItalicFont=cmunso.otf]
\newfontfamily\cyrillicfonttt{cmuntt.otf}[BoldFont=cmuntb.otf, ItalicFont=cmunit.otf, BoldItalicFont=cmuntx.otf]
\PolyglossiaSetup{russian}{hyphenmins={2,3}}

\usepackage[a4paper,left=8mm,right=8mm,top=7mm,bottom=11mm,footskip=5.5mm]{geometry}
\usepackage{xcolor}
\usepackage{multicol}
\usepackage{enumitem}
\usepackage{ragged2e}
\usepackage{fancyhdr}
\usepackage[colorlinks=true,linkcolor=hseblue,urlcolor=hseblue]{hyperref}

\definecolor{hseblue}{RGB}{16,45,105}
\definecolor{hseaccent}{RGB}{198,40,40}
\definecolor{hsegray}{RGB}{108,114,126}
\definecolor{hselight}{RGB}{240,243,248}

\hypersetup{pdftitle={Шпаргалка к семинару 6: Методы второго порядка, часть вторая}, pdfauthor={Евгений Баулин}}

\pagestyle{fancy}
\fancyhf{}
\renewcommand{\headrulewidth}{0pt}
\fancyfoot[L]{\scriptsize\color{hsegray}\textcopyright\ 2026 Евгений Баулин, \href{https://t.me/tarakan_tuc}{t.me/tarakan\_tuc}}
\fancyfoot[R]{\scriptsize\color{hsegray}Лицензия MIT}

\setlength{\parindent}{0pt}
\setlength{\parskip}{0pt}
\setlength{\columnsep}{4.5mm}
\setlength{\columnseprule}{0.3pt}
\renewcommand{\columnseprulecolor}{\color{hselight}}
\setlength{\emergencystretch}{2em}
\setlength{\RaggedRightRightskip}{0pt plus 4em}
\binoppenalty=2000
\relpenalty=500
\setlength{\abovedisplayskip}{2pt}
\setlength{\belowdisplayskip}{2pt}
\setlist[itemize]{leftmargin=7pt, labelsep=3pt, itemsep=1.2pt, topsep=1pt, parsep=0pt, partopsep=0pt,
  label={\color{hseaccent}\scriptsize\textbullet}}

\makeatletter
\newcommand{\cssec}[1]{\par\addvspace{3.5pt}\noindent
  \colorbox{hselight}{\parbox{\dimexpr\linewidth-2\fboxsep}{\RaggedRight\hyphenpenalty=10000\color{hseblue}\bfseries\strut #1}}%
  \par\nobreak\vspace{1pt}\@afterheading}
\makeatother
\clubpenalty=10000
\widowpenalty=10000

\begin{document}
\thispagestyle{fancy}

\noindent
\begin{minipage}[b]{0.6\textwidth}
  \RaggedRight\hyphenpenalty=10000
  {\color{hseaccent}\bfseries\small Шпаргалка к семинару 6}\\[1pt]
  {\color{hseblue}\bfseries\Large Методы второго\\ порядка, часть вторая}
\end{minipage}\hfill
\begin{minipage}[b]{0.39\textwidth}
  \raggedleft\footnotesize
  {\bfseries Методы оптимизации}\\
  НИУ ВШЭ, факультет компьютерных наук, 2026--2027\\
  Лектор курса: Андрей Игнатов\\
  Автор: Евгений Баулин, Telegram \href{https://t.me/tarakan_tuc}{@tarakan\_tuc}
\end{minipage}

\vspace{2pt}
{\color{hseblue}\rule{\textwidth}{0.8pt}}
\vspace{-4pt}

\fontsize{9}{10.6}\selectfont
\begin{multicols}{3}
\RaggedRight


% ---- gn
\cssec{Гаусс–Ньютон}
\begin{itemize}
  \item $f=\tfrac{1}{2}\|r\|^2$: $\nabla f=J^\top r$, $\nabla^2 f=J^\top J+\sum r_i\nabla^2 r_i$.
  \item GN: $J^\top Jp=-J^\top r$. Нужен полный столбцовый ранг $J$.
  \item Малые невязки возле $x^*$: квадратичная сходимость; большие — линейная.
\end{itemize}

% ---- lm
\cssec{Левенберг–Марквардт}
\begin{itemize}
  \item $(J^\top J+\lambda I)p=-J^\top r$; $\rho=\frac{f-f^+}{m(0)-m(p)}$.
  \item $\rho\ge3/4$: принять, $\lambda\mathrel{/}=10$. $\rho<1/4$: отвергнуть, $\lambda\mathrel{*}=10$.
  \item Малое $\lambda$: Гаусс–Ньютон. Большое: короткий градиентный шаг.
\end{itemize}

% ---- tr
\cssec{Доверительная область}
\begin{itemize}
  \item $\min m(p)$ при $\|p\|\le\Delta$: $(B+\lambda I)p=-g$ с $\lambda\ge0$.
  \item Dogleg: ломаная $0\to p_C\to p_N$; самая дальняя точка в шаре.
  \item $\rho>3/4$ на границе: $\Delta\leftarrow 2\Delta$. $\rho<1/4$: $\Delta\leftarrow\Delta/4$.
\end{itemize}

% ---- lbfgs
\cssec{L-BFGS}
\begin{itemize}
  \item Хранить $m$ пар $(s_i,y_i)$; $\gamma_k=s_{k-1}^\top y_{k-1}/(y_{k-1}^\top y_{k-1})$.
  \item Две петли: обратный проход $a_i,q$; масштаб $\gamma_k q$; прямой проход $b_i,r$.
  \item Одно направление: $4mn$ работы, $2mn$ памяти, без матрицы $n\times n$.
\end{itemize}

\end{multicols}
\end{document}
